\documentclass[12pt, titlepage]{article}

\usepackage{fullpage}
\usepackage[round]{natbib}
\usepackage{multirow}
\usepackage{booktabs}
\usepackage{tabularx}
\usepackage{graphicx}
\usepackage{float}
\usepackage{xr}
\usepackage{hyperref}
\hypersetup{
    colorlinks,
    citecolor=blue,
    filecolor=black,
    linkcolor=red,
    urlcolor=blue
}

  \externaldocument{../../SRS/SRS}

\input{../../Comments.text}
\input{../../Common.text}

\newcounter{acnum}
\newcommand{\actheacnum}{AC\theacnum}
\newcommand{\acref}[1]{AC\ref{#1}}

\newcounter{ucnum}
\newcommand{\uctheucnum}{UC\theucnum}
\newcommand{\uref}[1]{UC\ref{#1}}

\newcounter{mnum}
\newcommand{\mthemnum}{M\themnum}
\newcommand{\mref}[1]{M\ref{#1}}
\newcommand{\lcref}[1]{LC\ref{#1}}

\begin{document}

\title{Module Guide for \progname{}} 
\author{\authname}
\date{\today}

\maketitle

\pagenumbering{roman}

\section{Revision History}

\begin{tabularx}{\textwidth}{p{3cm}p{2cm}X}
\toprule {\bf Date} & {\bf Version} & {\bf Notes}\\
\midrule
Date 1 & 1.0 & Notes\\
Date 2 & 1.1 & Notes\\
\bottomrule
\end{tabularx}

\newpage

\section{Reference Material}

This section records information for easy reference. 

\subsection{Abbreviations and Acronyms}

\renewcommand{\arraystretch}{1.2}
\begin{tabular}{l l} 
  \toprule		
  \textbf{symbol} & \textbf{description}\\
  \midrule 
  AC & Anticipated Change\\
  DAG & Directed Acyclic Graph \\
  M & Module \\
  MG & Module Guide \\
  OS & Operating System \\
  R & Requirement\\
  SC & Scientific Computing \\
  SRS & Software Requirements Specification\\
  \progname & Explanation of program name\\
  UC & Unlikely Change \\
  \wss{etc.} & \wss{...}\\
  \bottomrule
\end{tabular}\\

\newpage

\tableofcontents

\listoftables

\listoffigures

\newpage

\pagenumbering{arabic}

\section{Introduction}

Decomposing a system into modules is a commonly accepted approach to developing software. A module is a work assignment for a programmer or programming team~\citep{ParnasEtAl1984}. We advocate a decomposition based on the principle of information hiding~\citep{Parnas1972a}. This principle supports design for change, because the ``secrets'' that each module hides represent likely future changes. Design for change is particularly valuable in software development, where modifications are frequent, especially during initial development as the solution space is explored.
The system described in this document is the Decentralized AI Inference Marketplace (DAIIM), a web-based platform that allows users to submit inference requests and receive aggregated results from participating inference providers. The system uses blockchain-based programs to manage marketplace operations, including request lifecycles, provider participation, prediction aggregation, and token-based settlement. Supporting off-chain components provide services such as user interaction, transaction management, request input handling, and communication with inference providers.
Our design follows the rules laid out by \citet{ParnasEtAl1984}, as follows:
\begin{itemize}
\item System details that are likely to change independently should be the secrets of separate modules.
\item Each data structure is implemented in only one module.
\item Any other program that requires information stored in a module's data structures must obtain it by calling access programs belonging to that module.
\end{itemize}
After completing the first stage of the design, the Software Requirements Specification (SRS), the Module Guide (MG) is developed~\citep{ParnasEtAl1984}. The MG specifies the modular structure of the system and is intended to allow both designers and maintainers to easily identify the parts of the software. The potential readers of this document are as follows:
\begin{itemize}
\item New project members: This document provides an overview of the system's modular structure, allowing new members to understand the overall design and quickly locate the modules relevant to their work.
\item Maintainers: The hierarchical structure of the module guide helps maintainers understand the responsibilities and relationships between modules when modifications are required. Maintainers should update the relevant sections of this document when changes are made to the system.
\item Designers: The module guide can be used to evaluate the consistency, feasibility, and flexibility of the system's decomposition. Designers can use it to verify that module responsibilities are appropriately separated and that anticipated changes can be accommodated without requiring unnecessary modifications to other modules.
\end{itemize}
The rest of the document is organized as follows. Section \ref{SecChange} lists the anticipated and unlikely changes to the system. Section \ref{SecMH} summarizes the module decomposition constructed according to the information-hiding principle. Section \ref{SecConnection} specifies the connections between the software requirements and the modules. Section \ref{SecMD} gives a detailed description of the modules. Section \ref{SecTM} includes two traceability matrices. One checks the completeness of the design against the requirements provided in the SRS. The other shows the relationship between anticipated changes and the modules. Section \ref{SecUse} describes the use relations between modules.

\section{Anticipated and Unlikely Changes} \label{SecChange}

This section identifies possible changes to DAIIM and classifies them according to their likelihood. Anticipated changes are listed in Section \ref{SecAchange}, while unlikely changes are listed in Section \ref{SecUchange}.

\subsection{Anticipated Changes} \label{SecAchange}

Anticipated changes are design or implementation decisions that may reasonably need to be modified as the system develops. These decisions should be hidden within the modules responsible for them so that a change can be made with minimal impact on the rest of the system.

\begin{description}

\item[\refstepcounter{acnum} \actheacnum \label{acInferenceModels}:] The types of AI inference tasks supported by DAIIM. The initial system focuses on regression tasks with numerical outputs, but support for additional inference types may be introduced in the future. The inference processing and aggregation components may need to be adapted to accommodate these additional types.

\item[\refstepcounter{acnum} \actheacnum \label{acAggregation}:] The algorithms used to aggregate provider predictions, update provider reputations, and calculate rewards. These mechanisms may be swapped out as the system is evaluated to improve prediction reliability, incentive fairness, and resistance to inaccurate or malicious providers.

\item[\refstepcounter{acnum} \actheacnum \label{acUserInterface}:] The layout and presentation of the web-based user interface. The interface may change based on user feedback and usability testing, particularly regarding task creation, provider information, transaction status, and the presentation of aggregated results and prediction quality.

\item[\refstepcounter{acnum} \actheacnum \label{acInfrastructure}:] The implementation of off-chain supporting services and their communication mechanisms. The backend, database, and provider SDK may need to be updated or replaced as the system evolves to improve scalability, maintainability, or compatibility with external services.


\end{description}

\subsection{Unlikely Changes} \label{SecUchange}

Although the system should accommodate anticipated changes, certain fundamental design decisions are not expected to change during the scope of the project. These decisions establish the core architecture and purpose of DAIIM, and changing them would likely require modifications across multiple parts of the system.

\begin{description}

\item[\refstepcounter{ucnum} \uctheucnum \label{ucCorePurpose}:] The fundamental purpose of DAIIM as a decentralized marketplace for requesting AI inferences from multiple independent providers and combining their predictions into an aggregated result. The system is not intended to become a centralized inference service controlled by a single provider.

\item[\refstepcounter{ucnum} \uctheucnum \label{ucBlockchainArchitecture}:] The use of a public blockchain as an essential part of the system architecture. Marketplace operations such as request management, provider registration, prediction aggregation, reputation updates, and token settlement will continue to rely on on-chain execution and records rather than being controlled entirely by the off-chain backend.

\item[\refstepcounter{ucnum} \uctheucnum \label{ucWebPlatform}:] The system's intended form as a web-based application. Task requesters and inference providers will continue to interact with DAIIM through a web interface or its provider-facing SDK, rather than requiring DAIIM to become a standalone desktop or mobile application.

\end{description}


\section{Module Hierarchy} \label{SecMH}

This section provides an overview of the module design.
Modules are summarized in a hierarchy decomposed by secrets in Table \ref{TblMH}.
The modules listed below, which are leaves in the hierarchy tree, are the modules that will actually be implemented.

\begin{description}

\item [\refstepcounter{mnum} \mthemnum \label{mWallet}:] Wallet Manager

\item [\refstepcounter{mnum} \mthemnum \label{mChain}:] Chain Observer

\item [\refstepcounter{mnum} \mthemnum \label{mInput}:] Request Input Manager

\item [\refstepcounter{mnum} \mthemnum \label{mUI}:] User Interface

\item [\refstepcounter{mnum} \mthemnum \label{mUserInfo}:] User Information

\item [\refstepcounter{mnum} \mthemnum \label{mProviderSetup}:] Provider Setup

\item [\refstepcounter{mnum} \mthemnum \label{mListener}:] Request Listener

\item [\refstepcounter{mnum} \mthemnum \label{mParticipation}:] Participation Client

\item [\refstepcounter{mnum} \mthemnum \label{mRegistry}:] Provider Registry

\item [\refstepcounter{mnum} \mthemnum \label{mLifecycle}:] Request Lifecycle

\item [\refstepcounter{mnum} \mthemnum \label{mSettlement}:] Settlement

\item [\refstepcounter{mnum} \mthemnum \label{mToken}:] Token Manager

\item [\refstepcounter{mnum} \mthemnum \label{mTransaction}:] Transaction Manager

\end{description}

\begin{table}[h!]
\centering
\caption{Module Hierarchy}
\label{TblMH}
\begin{tabular}{p{0.3\textwidth} p{0.6\textwidth}}
\toprule
\textbf{Level 1} & \textbf{Level 2}\\
\midrule

\multirow{5}{0.3\textwidth}{Behaviour-Hiding Module}
& \mref{mWallet}: Wallet Manager \\
& \mref{mChain}: Chain Observer \\
& \mref{mInput}: Request Input Manager \\
& \mref{mUI}: User Interface \\
& \mref{mUserInfo}: User Information \\
& \mref{mProviderSetup}: Provider Setup \\
& \mref{mListener}: Request Listener \\
& \mref{mParticipation}: Participation Client \\
\midrule

\multirow{5}{0.3\textwidth}{Software Decision Module}
& \mref{mRegistry}: Provider Registry \\
& \mref{mLifecycle}: Request Lifecycle \\
& \mref{mSettlement}: Settlement \\
& \mref{mToken}: Token Manager \\
& \mref{mTransaction}: Transaction Manager \\

\bottomrule
\end{tabular}
\end{table}

\section{Connection Between Requirements and Design} \label{SecConnection}

The design of DAIIM is intended to satisfy the functional requirements developed in the SRS. The following key design decisions were made to address specific requirements:

\begin{itemize}

    \item \textbf{Wallet Manager (\mref{mWallet}):} To satisfy FR\ref{FR_WalletAuthentication} and FR\ref{FR_TransactionAuthorization}, wallet authentication and transaction signing are handled through a dedicated wallet module. This module verifies ownership of wallet addresses and provides the interface for users to authorize transactions without exposing their private keys to DAIIM.

    \item \textbf{Chain Observer (\mref{mChain}):} To satisfy FR\ref{FR_OnChainRecords}, FR\ref{FR_TransactionStatus}, and FR\ref{FR_TransactionVerification}, a dedicated module monitors blockchain activity and retrieves transaction and account information. This separates reading and interpreting blockchain data from the modules responsible for initiating transactions and managing marketplace operations.

    \item \textbf{Request Input Manager (\mref{mInput}):} To satisfy FR\ref{FR_RequestInputDistribution}, FR\ref{FR_InputAccessRestriction}, FR\ref{FR_InputNotOnChain}, and FR\ref{FR_InputDeletion}, a dedicated module manages the secure storage, encryption, access, and deletion of inference request inputs. Input data is kept off-chain and made available only to the requester and providers that accepted the corresponding request.

    \item \textbf{User Interface (\mref{mUI}):} To satisfy FR\ref{FR_PublicAccess}, FR\ref{FR_TaskBrowsing}, FR\ref{FR_InferenceTaskCreation}, and FR\ref{FR_FinalResultDelivery}, the web interface provides users with access to marketplace information, task creation, inference requests, and aggregated results. It also presents transaction information and marketplace activity in a way that does not require users to interact directly with the blockchain.

    \item \textbf{User Information (\mref{mUserInfo}):} To support FR\ref{FR_WalletAuthentication} and FR\ref{FR_ProviderDashboard}, a dedicated module manages user profile information associated with verified wallet addresses. This allows user-specific information, such as provider activity and preferences, to be maintained separately from the underlying blockchain records.

    \item \textbf{Provider Setup (\mref{mProviderSetup}):} To support FR\ref{FR_InferenceTaskRegistrationStaking} and FR\ref{FR_ProviderDashboard}, a dedicated module handles the provider's initial connection to the blockchain and retrieves the provider's registration information for available tasks. This separates provider initialization and configuration from the ongoing handling of inference requests.

    \item \textbf{Request Listener (\mref{mListener}):} To satisfy FR\ref{FR_RequestBroadcast} and FR\ref{FR_AcceptancePeriodStart}, a dedicated module monitors inference requests that are relevant to a provider's registered tasks. It identifies new requests and communicates their details, payment information, and applicable deadlines to the provider.

    \item \textbf{Participation Client (\mref{mParticipation}):} To satisfy FR\ref{FR_RequestInputDistribution}, FR\ref{FR_PredictionSubmission}, FR\ref{FR_OnePredictionPerProvider}, and FR\ref{FR_PredictionValidation}, a dedicated client module handles the provider's participation in inference requests. It retrieves authorized input data, generates the prediction commitment using a cryptographically secure salt, and submits commitments and reveals directly to the on-chain program.

    \item \textbf{Provider Registry (\mref{mRegistry}):} To satisfy FR\ref{FR_InferenceTaskRegistrationStaking}, FR\ref{FR_ProviderWillingUnregistration}, and FR\ref{FR_ProviderMinimumStakeUnregistration}, a dedicated module manages provider registration for individual inference tasks. It maintains provider eligibility, registration status, and task-specific staking requirements so that only eligible providers can participate in requests.

    \item \textbf{Request Lifecycle (\mref{mLifecycle}):} To satisfy FR\ref{FR_RequestSubmission} through FR\ref{FR_MinimumPredictions}, a dedicated module manages the state and progression of inference requests. The request lifecycle is implemented as a state machine that controls the acceptance, commit, and reveal phases, enforces their deadlines, validates provider participation, and determines whether a request can proceed or must be cancelled.

    \item \textbf{Settlement (\mref{mSettlement}):} To satisfy FR\ref{FR_PredictionAggregation}, FR\ref{FR_UpdateProviderReputations}, FR\ref{FR_ReputationRange}, and FR\ref{FR_DistributeRewards}, a dedicated module handles the finalization of inference requests. It calculates the aggregated result using the reputation-weighted median algorithm, updates provider reputations using the confidence-based reputation algorithm, and determines the distribution of rewards based on the resulting provider reputations.

    \item \textbf{Token Manager (\mref{mToken}):} To satisfy FR\ref{FR_RequestPayment}, FR\ref{FR_StakeReservation}, FR\ref{FR_StakeReservationRelease}, and FR\ref{FR_InferenceProviderSlashing}, a dedicated module manages token-related operations. These include transferring payments into escrow, managing provider stakes, releasing reserved funds, distributing tokens, and applying penalties such as slashing or burning when required.

    \item \textbf{Transaction Manager (\mref{mTransaction}):} To satisfy FR\ref{FR_TransactionPreview}, FR\ref{FR_TransactionStatus}, and FR\ref{FR_TransactionVerification}, a dedicated module coordinates the preparation and submission of user-authorized transactions. It validates transaction parameters, checks available balances, simulates transactions before submission, and tracks their status to prevent duplicate submissions or incorrect reporting of transaction outcomes.

\end{itemize}


\section{Module Decomposition} \label{SecMD}

Modules are decomposed according to the principle of ``information hiding''
proposed by \citet{ParnasEtAl1984}. The \emph{Secrets} field in a module
decomposition is a brief statement of the design decision hidden by the
module. The \emph{Services} field specifies \emph{what} the module will do
without documenting \emph{how} to do it. For each module, a suggestion for the
implementing software is given under the \emph{Implemented By} title. If the
entry is \emph{OS}, this means that the module is provided by the operating
system or by standard programming language libraries.  \emph{\progname{}} means the
module will be implemented by the \progname{} software.

Only the leaf modules in the hierarchy have to be implemented. If a dash
(\emph{--}) is shown, this means that the module is not a leaf and will not have
to be implemented.

\subsection{Hardware Hiding Modules (\mref{mHH})}

\begin{description}
\item[Secrets:]The data structure and algorithm used to implement the virtual
  hardware.
\item[Services:]Serves as a virtual hardware used by the rest of the
  system. This module provides the interface between the hardware and the
  software. So, the system can use it to display outputs or to accept inputs.
\item[Implemented By:] OS
\end{description}

\subsection{Behaviour-Hiding Module}

\begin{description}
\item[Secrets:]The contents of the required behaviours.
\item[Services:]Includes programs that provide externally visible behaviour of
  the system as specified in the software requirements specification (SRS)
  documents. This module serves as a communication layer between the
  hardware-hiding module and the software decision module. The programs in this
  module will need to change if there are changes in the SRS.
\item[Implemented By:] --
\end{description}

\subsubsection{Input Format Module (\mref{mInput})}

\begin{description}
\item[Secrets:]The format and structure of the input data.
\item[Services:]Converts the input data into the data structure used by the
  input parameters module.
\item[Implemented By:] [Your Program Name Here]
\item[Type of Module:] [Record, Library, Abstract Object, or Abstract Data Type]
  [Information to include for leaf modules in the decomposition by secrets tree.]
\end{description}

\subsubsection{Etc.}


\subsection{Software Decision Module}

\begin{description}
\item[Secrets:] The design decision based on mathematical theorems, physical
  facts, or programming considerations. The secrets of this module are
  \emph{not} described in the SRS.
\item[Services:] Includes data structure and algorithms used in the system that
  do not provide direct interaction with the user. 
  % Changes in these modules are more likely to be motivated by a desire to
  % improve performance than by externally imposed changes.
\item[Implemented By:] --
\end{description}

\subsubsection{Etc.}

\section{Traceability Matrix} \label{SecTM}

This section shows the traceability between the system modules and the
requirements, as well as between the modules and the likely and unlikely
changes anticipated for DAIIM. These matrices demonstrate which modules
are responsible for satisfying the requirements and which modules may
need to be modified if the system evolves.

\begin{table}[H]
\centering
\caption{Trace Between Requirements and Modules}
\label{TblRT}
\begin{tabularx}{\textwidth}{|X|X|}
\hline
\textbf{Requirement(s)} & \textbf{Modules}\\
\hline

FR\ref{FR_WalletAuthentication}, FR\ref{FR_EmbeddedWalletCreation},
FR\ref{FR_TransactionAuthorization}
& \mref{mWallet}, \mref{mUserInfo}, \mref{mUI}, \mref{mTransaction} \\
\hline

FR\ref{FR_PublicAccess}, FR\ref{FR_TaskBrowsing}
& \mref{mUI}, \mref{mRegistry}, \mref{mUserInfo} \\
\hline

FR\ref{FR_InferenceTaskCreation}--FR\ref{FR_TaskAcceptancePeriod}
& \mref{mUI}, \mref{mLifecycle}, \mref{mTransaction}, \mref{mToken} \\
\hline

FR\ref{FR_InferenceTaskRegistration},
FR\ref{FR_InferenceTaskRegistrationStaking},
FR\ref{FR_ProviderWillingUnregistration},
FR\ref{FR_ProviderMinimumStakeUnregistration}
& \mref{mRegistry}, \mref{mToken}, \mref{mTransaction},
\mref{mWallet}, \mref{mProviderSetup} \\
\hline

FR\ref{FR_ProviderDashboard}
& \mref{mUI}, \mref{mUserInfo}, \mref{mRegistry},
\mref{mChain} \\
\hline

FR\ref{FR_RequestSubmission}, FR\ref{FR_RequestValidation},
FR\ref{FR_RequestPayment}, FR\ref{FR_RequestBroadcast}
& \mref{mUI}, \mref{mLifecycle}, \mref{mInput},
\mref{mToken}, \mref{mTransaction}, \mref{mListener} \\
\hline

FR\ref{FR_AcceptancePeriodStart}--FR\ref{FR_MaximumAcceptance},
FR\ref{FR_MinimumAcceptance}
& \mref{mLifecycle}, \mref{mRegistry}, \mref{mListener},
\mref{mToken} \\
\hline

FR\ref{FR_RequestInputDistribution},
FR\ref{FR_InputAccessRestriction},
FR\ref{FR_InputNotOnChain},
FR\ref{FR_InputDeletion}
& \mref{mInput}, \mref{mParticipation}, \mref{mLifecycle} \\
\hline

FR\ref{FR_PredictionPhaseStart},
FR\ref{FR_ProtocolCommitPeriod}--FR\ref{FR_ProtocolRevealPeriod},
FR\ref{FR_PredictionSubmission}--FR\ref{FR_PredictionValidation}
& \mref{mLifecycle}, \mref{mParticipation},
\mref{mInput}, \mref{mListener} \\
\hline

FR\ref{FR_PredictionPhaseEnd},
FR\ref{FR_MinimumPredictions},
FR\ref{FR_StakeReservation}
& \mref{mLifecycle}, \mref{mRegistry}, \mref{mToken},
\mref{mParticipation} \\
\hline

FR\ref{FR_PredictionAggregation},
FR\ref{FR_FinalResultDelivery},
FR\ref{FR_ResultSpread},
FR\ref{FR_RequestHistory}
& \mref{mSettlement}, \mref{mLifecycle}, \mref{mUI},
\mref{mChain} \\
\hline

FR\ref{FR_UpdateProviderReputations},
FR\ref{FR_ReputationRange},
FR\ref{FR_DistributeRewards}
& \mref{mSettlement}, \mref{mRegistry}, \mref{mToken} \\
\hline

FR\ref{FR_RewardTransferFailure}--FR\ref{FR_RequestCompletion},
FR\ref{FR_CancellationCompensation}
& \mref{mSettlement}, \mref{mLifecycle}, \mref{mToken} \\
\hline

FR\ref{FR_InferenceProviderSlashing}--FR\ref{FR_ProviderForcedUnregistration}
& \mref{mRegistry}, \mref{mSettlement}, \mref{mToken},
\mref{mLifecycle} \\
\hline

FR\ref{FR_TransactionPreview},
FR\ref{FR_TransactionStatus},
FR\ref{FR_TransactionVerification},
FR\ref{FR_OnChainRecords}
& \mref{mTransaction}, \mref{mWallet}, \mref{mChain},
\mref{mUI}, \mref{mToken} \\
\hline

SA\ref{SA_NoDuplicateSubmission}--SA\ref{SA_FailureRetry},
SA\ref{SA_TransactionSimulation},
SA\ref{SA_ChargeEqualsAuthorized},
SA\ref{SA_RewardValidation}
& \mref{mTransaction}, \mref{mWallet}, \mref{mToken},
\mref{mUI}, \mref{mSettlement} \\
\hline

NFR\ref{NFR_DataEncryption},
NFR\ref{NFR_InputDataAccess}
& \mref{mInput}, \mref{mParticipation} \\
\hline

NFR\ref{NFR_PredictionsHidden},
NFR\ref{NFR_CommitSalt}
& \mref{mParticipation}, \mref{mLifecycle} \\
\hline

NFR\ref{NFR_ReputationWeightedAggregate},
NFR\ref{NFR_AggregationArithmeticErrors}
& \mref{mSettlement}, \mref{mToken} \\
\hline

NFR\ref{NFR_FinalState},
NFR\ref{NFR_ConcurrentRequests}
& \mref{mLifecycle}, \mref{mSettlement}, \mref{mRegistry},
\mref{mToken} \\
\hline

NFR\ref{NFR_ProgressStatus},
NFR\ref{NFR_TransactionConfirmation},
NFR\ref{NFR_ErrorVisibility}
& \mref{mUI}, \mref{mTransaction}, \mref{mChain} \\
\hline

NFR\ref{NFR_TwoClickAccess},
NFR\ref{NFR_PublicPagesPrivacy},
NFR\ref{NFR_Learnability},
NFR\ref{NFR_BrowserSupport}
& \mref{mUI}, \mref{mWallet}, \mref{mUserInfo} \\
\hline

\end{tabularx}
\end{table}


\begin{table}[H]
\centering
\caption{Trace Between Likely Changes and Modules}
\label{TblLCT}
\begin{tabularx}{\textwidth}{|X|X|}
\hline
\textbf{Likely Change (LC)} & \textbf{Modules}\\
\hline

\lcref{LC_inferenceModels}
& \mref{mProviderSetup}, \mref{mParticipation},
\mref{mLifecycle}, \mref{mSettlement} \\
\hline

\lcref{LC_providerParticipation}
& \mref{mRegistry}, \mref{mLifecycle}, \mref{mSettlement},
\mref{mToken}, \mref{mListener} \\
\hline

\lcref{LC_aggregation}
& \mref{mSettlement}, \mref{mRegistry}, \mref{mToken} \\
\hline

\lcref{LC_randomInvitations}
& \mref{mListener}, \mref{mRegistry}, \mref{mLifecycle} \\
\hline

\lcref{LC_userRequirements}
& \mref{mUI}, \mref{mUserInfo} \\
\hline

\end{tabularx}
\end{table}


\begin{table}[H]
\centering
\caption{Trace Between Unlikely Changes and Modules}
\label{TblUCT}
\begin{tabularx}{\textwidth}{|X|X|}
\hline
\textbf{Unlikely Change (LC)} & \textbf{Modules}\\
\hline

\lcref{LC_corePurpose}
& \mref{mUI}, \mref{mRegistry}, \mref{mLifecycle},
\mref{mSettlement}, \mref{mParticipation} \\
\hline

\lcref{LC_softwarePlatform}
& \mref{mUI}, \mref{mWallet}, \mref{mProviderSetup} \\
\hline

\end{tabularx}
\end{table}

\section{Use Hierarchy Between Modules} \label{SecUse}

In this section, the uses hierarchy between modules is provided.
A module A uses a module B if the correct execution of A depends upon the availability of a correct implementation of B. The arrows in the uses hierarchy are directed from the module that requires a service to the module that provides that service.
For DAIIM, the uses hierarchy is organized around the User Interface, the provider-side SDK modules, the supporting off-chain modules, and the core on-chain modules.
The User Interface (\mref{mUI}) depends on the Wallet Manager (\mref{mWallet}), User Information (\mref{mUserInfo}), Transaction Manager (\mref{mTransaction}), and Chain Observer (\mref{mChain}) to support wallet authentication, user information retrieval, transaction submission, and the display of blockchain state.
On the provider side, the Provider Setup module (\mref{mProviderSetup}) uses the Wallet Manager (\mref{mWallet}), User Information (\mref{mUserInfo}), and Chain Observer (\mref{mChain}) to establish provider identity and retrieve registration information. The Request Listener (\mref{mListener}) uses the Chain Observer (\mref{mChain}) to monitor relevant requests and blockchain events. The Participation Client (\mref{mParticipation}) uses the Provider Setup module (\mref{mProviderSetup}), Request Input Manager (\mref{mInput}), and Transaction Manager (\mref{mTransaction}) to retrieve authorized request inputs and participate in inference tasks.
The Transaction Manager (\mref{mTransaction}) uses the Wallet Manager (\mref{mWallet}) and Chain Observer (\mref{mChain}) to prepare, authorize, and verify transactions.
The core on-chain modules provide the functionality required to operate the marketplace. The Request Lifecycle module (\mref{mLifecycle}) uses the Provider Registry (\mref{mRegistry}) and Token Manager (\mref{mToken}) to manage provider eligibility and request-related token operations. The Settlement module (\mref{mSettlement}) uses the Request Lifecycle module (\mref{mLifecycle}), Provider Registry (\mref{mRegistry}), and Token Manager (\mref{mToken}) to calculate final results, update provider reputations, and distribute rewards or penalties.
Figure \ref{FigUH} illustrates the uses relationships between the modules. The resulting hierarchy is a directed acyclic graph (DAG), with higher-level modules depending on lower-level modules that provide the services they require.

\begin{figure}[H]
\centering
%\includegraphics[width=0.7\textwidth]{UsesHierarchy.png}
\caption{Use hierarchy among modules}
\label{FigUH}
\end{figure}

%\section*{References}

\section{User Interfaces}

\wss{Design of user interface for software and hardware.  Attach an appendix if
needed. Drawings, Sketches, Figma}

\section{Design of Communication Protocols}

\wss{If there is communication between hardware and software, it may be necessary 
to explain the design of the communication protocol.  This section should be filled 
in with NA if it is not appropriate for a given project.}

\section{Database Design}

\wss{If a database is a significant component of the project, the database
design should be summarized here.  This section should be filled in with NA if
it is not appropriate for a given project.}

\section{Timeline}

\wss{Schedule of tasks and who is responsible}

\wss{You can point to GitHub if this information is included there}

\section{Generative AI Use Disclosure}

\wss{ChatGPT (version 5) was used to brainstorm ideas for the template for this disclose section.}

\subsection{AI Tools Used}

\wss{Give the name of the AI tool(s) used and the version.}

\subsection{How and Where Used}

\wss{Explain what the tool(s) were used for.  Examples include: brainstorming,
explaining a technical concept, reviewing the document, generating or modifying
text, generating or modifying diagrams, identifying errors or omissions or
inconsistencies.}

\wss{For each of the uses, identify which parts of the document were affected.}

\wss{You do not need to report every interaction, but you should disclose the 
uses that materially contributed to your work}

\subsection{Verification of AI-Generated Content}

\wss{\begin{itemize}
    \item Describe how the team checked AI-generated or AI-assisted material
    before including it in the document.
    \item In particular, verify factual claims, technical assertions,
    requirements, calculations, references, and other information for which an
    AI system may produce incorrect or fabricated results.
    \item \textbf{The team is responsible for the accuracy, completeness,
consistency, and appropriateness of the submitted document, regardless of
whether material was generated by a human or an AI system.}
\end{itemize}}

\newpage{}

\bibliographystyle {plainnat}
\bibliography{../../../refs/References}

\end{document}